\chapter{Einführung}

Asynchrone und digitale Kommunikationstechniken sind von der heutigen Welt nicht
mehr wegzudenken, welche für Einsatzgebiete wie der Kommunikation zwischen Freunden
bis hin zum Nutzen in großen Konzernen ragen. Durch die Internationalität in solchen
großen Unternehmen werden Kommunikations-Möglichkeiten benötigt, welche den Austausch
und das Festhalten von Informationen über Zeit und Ländergrenzen hinweg ermöglichen.

Auf dem Markt für solche Kommunikationsmittel stehen viele verschiedene Produkte
zur Verfügung. Eines hierfür, welches eine große Beliebtheit erfährt und weitere
Produkte maßgeblich beeinflusst hat ist \emph{Slack}.~\cite{slack} Dieses ist
im Kern ein Chatclient, welche mit verschiedenen Textkanälen in einem Team für
die Kollaboration von Teammitgliedern arbeitet. So kann eine Organisation erstellt werden,
welche alle Nutzer eines Unternehmens beingehören. In dieser Organisation stehen viele einzelne Textkanäle bereit,
in welchen sie die Organisationsmitglieder austauschen können. Neben der Kommunikation in großen Kanälen
gibt es auch die Möglichkeit als Person mit einzelnen Personen des Teams sich über
private Nachrichten auszutauschen. Durch die Beliebtheit von \emph{Slack} in Unternehmen sind
auch weitere Chatclients, welche ein ähnliches Prinzip verfolgen, dies jedoch für andere Personengruppen anbieten, entstanden.
Ein weiteres beliebtes Produkt wäre \emph{Discord}~\cite{discord}, welches eher eine hohe private Nutzung erfährt.

So wurde in diesem Projekt ein ähnlicher Chatclient, namens \anf{Chatter}, wie die bisher
genannten Produkte, erstellt. Dieser verfolgt das Ziel eine Selbstorganisation von Nutzern
in verschiedene einzelne Teams zu erlauben und so einen effizienten Kommunikationsweg
für \zb{} das Arbeiten in Gruppen zu ermöglichen. Jeder Nutzer kann dafür selbst eigene
Teams erstellen und weitere Nutzer zu diesen einladen. In den Teams können über mehrere Kanäle
Textnachrichten ausgetauscht werden.
